\begin{corollary}[Integral Mordell--Weil lattice]\hypertarget{tag-0309}{}\label{push:src:B.3.5}\label{push:B35}
\StatementRecord{push:src:B.3.5}{proved}{0309}{\SeriesRef{push:src:B.3.4}}

In the notation of \ref{push:src:B.3.4},
\[
 \mathscr E(\mathbb C(C))
   =\mathbb ZP\oplus\mathbb ZR\oplus\mathbb ZR',
 \qquad \det\operatorname{NS}(\mathscr E)=-32.
\]
\end{corollary}


\begin{proof}
\begin{enumerate}[label=\arabic*.,leftmargin=*]
\item The affine coordinate ring of $C$ is
$\mathbb C[a,b]/(b^3-a^4-1)$, with
unique representatives of degree at most
two in $b$. The pole orders $3i+4j$ of
$a^ib^j$, $0\le j\le2$, are distinct,
so cancellation cannot create functions
with pole order one or two at infinity.
Consequently
\[
 H^0(C,\mathcal O_C(2\infty))=\mathbb C,
 \qquad
 H^0(C,\mathcal O_C(3\infty))
       =\mathbb C\oplus\mathbb Ca.
\]
A section disjoint from $O$ therefore
has $X=x_0$, $Y=y_0+y_1a$ for constants
$x_0,y_0,y_1$. Substitution in the
Weierstrass equation forces $x_0=y_0=0$
and $y_1^2=2$. Thus the only nonzero
sections with intersection zero with
$O$ are $P,-P$.
Every other nonzero section has
intersection at least one with $O$,
and hence height at least $4-3=1$.

\item \ref{push:src:B.3.4} proves that every section has
integral coefficient along $P$ in the
rational height decomposition. The
remaining lattice can have index only
one, two or four over $\mathbb ZR+\mathbb ZR'$.
Any proper such overlattice contains a
nonzero element of order two modulo
this sublattice. After subtracting an
integral combination of $R,R'$ it is
one of $R/2$, $R'/2$, $(R+R')/2$.
Each has height $1/3$, contradicting
the preceding lower bound. Thus the
index is one. The determinant already
computed in \ref{push:src:B.3.4} is the full
N\'eron--Severi determinant.

\end{enumerate}
\end{proof}
